\title{FlashAttention-3: Fast and Accurate Attention with Asynchrony and Low-precision}

\begin{abstract}
Attention mechanisms form a foundational component in Transformer-based models, yet they present significant computational bottlenecks, particularly in scaling large language models and enabling efficient long-context processing. \textsc{FlashAttention} introduced a technique for enhancing attention speed on GPUs by significantly reducing memory access operations. Despite this, the method does not fully capitalize on recent hardware innovations, as evidenced by \textsc{FlashAttention-2}'s utilization rate of only 35\% on the H100 GPU. In this work, we introduce three primary strategies to further accelerate attention on Hopper architecture GPUs: first, leveraging the asynchronous capabilities of Tensor Cores and TMA for (1) overlapping computation with data transfer using warp specialization and (2) fusing block-level matrix multiplication with softmax operations; second, employing block quantization and incoherent processing made feasible by native FP8 support. Through these approaches, our algorithm, \textsc{FlashAttention-3}, delivers a performance gain of 1.5-2.0$\times$ on H100 GPUs, achieving up to 740 TFLOPs/s in FP16 mode (corresponding to 75\% utilization) and nearly 1.2 PFLOPs/s using FP8 precision. Furthermore, we demonstrate that FP8 \textsc{FlashAttention-3} reduces numerical error by a factor of 2.6$\times$ compared to a standard FP8 attention baseline.
\end{abstract}

\section{Introduction}
\label{sec:intro}

Within the Transformer architecture~\citep{vaswani2017attention}, the attention mechanism serves as the main computational bottleneck, primarily due to the fact that calculating self-attention scores for queries and keys increases quadratically with respect to sequence length. Enhancing attention mechanisms to accommodate extended contexts has the potential to introduce new functionalities, such as processing and reasoning across multiple lengthy documents~\citep{guo2021longt5,shaham2022scrolls,peng2023yarn} or navigating through extensive codebases~\citep{roziere2023code, li2023starcoder}. It also enables novel modalities, including high-resolution visual data~\citep{chen2022scaling}, audio streams~\citep{gulati2020conformer}, and video~\citep{ho2022video}, as well as new use cases like managing long sequences of user interactions~\citep{sun2019bert4rec} and facilitating agent workflows that span substantial horizons~\citep{yao2022react}. These opportunities have spurred a great deal of research into accelerating attention in the context of lengthy inputs, either by introducing various approximations~\citep{katharopoulos2020transformers,choromanski2020rethinking,
tay2020efficient}, improving software implementations (\citep{rabe2021self,dao2022flashattention,kwon2023efficient}), or exploring entirely new architectural paradigms~\citep{peng2023rwkv,sun2023retentive,gu2023mamba}.

This study extends the foundational contributions of \citet{dao2022flashattention}, who devised precise attention computation methods that explicitly account for the GPU's architectural and operational features within their overall algorithmic framework. The initial proposal, \textsc{FlashAttention}, in \citep{dao2022flashattention} by Dao et al., introduced an innovative tiling method that enables the parallel execution of attention mechanisms while eliminating the need for slow intermediate transfers to global memory by consolidating all attention steps into a single GPU kernel. In a subsequent advancement, \citet{dao2023flashattention2} reformulated the approach into \textsc{FlashAttention-2}, extending the parallelization to the sequence length dimension and reordering the forward pass to process the key and value matrices in block segments. This redesign led to enhanced workload distribution and higher GPU occupancy. Despite these improvements, our analysis reveals that \textsc{FlashAttention-2} still underperforms compared to highly optimized matrix-multiplication (GEMM) kernels on recent GPUs, attaining utilization rates of only 35\% versus 80–90\% for GEMM on the Hopper H100 platform. Part of this inefficiency appears to stem from differences in implementation, such as the absence of Hopper-optimized instructions for the Tensor Cores, instead relying on Ampere-generation instructions. Recent studies, including ThunkerKitten~\citep{spector2024thunder} and cuDNN 9~\citep{cudnn9}, demonstrate that leveraging Hopper-specific instruction sets and employing tile-based abstractions can both accelerate attention operations and lead to a more streamlined implementation.

At a deeper level, the algorithm underlying \textsc{FlashAttention-2} operates within a basic synchronous paradigm and does not deliberately incorporate asynchronous processing or reduced-precision computation. Asynchronous execution emerges due to hardware innovations that enhance the performance of critical machine learning operations: for example, certain specialized components are tasked with executing matrix multiplications (such as Tensor Cores), or with managing memory transfers (via Tensor Memory Accelerator -- TMA), functioning independently of other CUDA cores responsible for general logical, integer, and floating-point tasks. Meanwhile, employing lower-precision formats like FP8 (introduced in Hopper) or FP4 (debuted in Blackwell)—an evolution from earlier use of FP16 (in Pascal, 2017) and BF16 (in Ampere, 2020)—has consistently enabled increased throughput, offering two to four times greater performance for comparable power and silicon area. The hardware capabilities provided by Hopper along these lines are surveyed in \cref{subsec:hardware}.

A central technical problem is reformulating \textsc{FlashAttention-2} to capitalize on these hardware advancements. Leveraging asynchrony involves orchestrating the overlap between matrix multiplication and softmax computations, even though their execution is interdependent. Utilizing low-precision arithmetic necessitates strategies to mitigate quantization errors, with particular attention to handling outlier activations that can arise in LLMs~\citep{dettmers2208llm, sun2024massive}.

In pursuit of this goal, we introduce \textsc{FlashAttention-3}, which integrates and advances three novel concepts to enhance efficiency on modern GPU platforms:\footnote{Our findings are presented using NVIDIA’s Hopper architecture as a reference, but the proposed algorithm remains applicable to any GPU system equipped with advanced asynchronous operations and support for low-precision arithmetic.}

\begin{enumerate}

\item \textbf{Producer-Consumer asynchrony:} We introduce a warp-specialized approach to software pipelining that leverages the asynchronous behavior of data transfers and Tensor Core operations. This is achieved by assigning data producers and consumers to distinct warps, which broadens the algorithm’s capacity to mask both memory and instruction dispatch latencies.

\item \textbf{Hiding softmax under asynchronous block-wise GEMMs:} We conceal the relatively low-throughput, non-GEMM tasks associated with the softmax operation—such as floating-point fused multiply-adds and exponentiations—by executing them concurrently with asynchronous WGMMA instructions for GEMM. To facilitate this overlap, we restructure the \textsc{FlashAttention-2} algorithm, eliminating certain sequential dependencies between the softmax and the GEMMs. For instance, in the two-stage variant of our method, as softmax is performed on one segment of the score matrix, the WGMMA operates asynchronously to prepare the next segment.

\item \textbf{Hardware-accelerated low-precision GEMM:} We modify the forward pass algorithm to target FP8 Tensor Cores during GEMM operations, effectively nearly doubling the observed TFLOPs/s. Accomplishing this necessitates reconciling the differing memory layout expectations for FP32 accumulators and FP8 operand matrices used by WGMMA. We apply strategies such as block quantization and incoherent processing to counteract the accuracy degradation that can occur when transitioning to FP8 computation.
\end{enumerate}

To empirically assess the effectiveness of our approach, we conducted benchmarks of \textsc{FlashAttention-3} on the H100 SXM5 GPU across various parameter settings. Our results demonstrate that: (1) in the forward pass, the FP16 configuration delivers a 1.5–2.0$\times$ speedup compared to \textsc{FlashAttention-2}, achieving up to 740 TFLOPs/s, and in the backward pass, it attains a 1.5–1.75$\times$ acceleration; (2) with FP8, throughput reaches nearly 1.2 PFLOPs/s; (3) for longer sequence lengths, FP16 shows clear advantages, while FP8 remains competitive with the most advanced attention implementation in NVIDIA's cuDNN library\footnote{Specifically, for head dimension 64, FP8 \textsc{FlashAttention-3} leads, whereas for head dimensions 128 and 256, it matches in cases without causal masking and trails slightly with causal masking.}. Furthermore, we confirm that FP16 \textsc{FlashAttention-3} exhibits the same level of numerical error as \textsc{FlashAttention-2}, and provides improved precision relative to the standard attention implementation, as critical intermediate values such as softmax rescaling are maintained in FP32. Additionally, in scenarios involving outlier features, FP8 \textsc{FlashAttention-3} equipped with block quantization and incoherent processing proves to be 2.6$\times$ more accurate than standard attention methods relying on per-tensor quantization.

We have released \textsc{FlashAttention-3} under a permissive open-source license\footnote{\textsc{FlashAttention-3} is available at \texttt{https://github.com/Dao-AILab/flash-attention}}. Our future plans include incorporating it into both the PyTorch and Hugging Face libraries to maximize its accessibility and utility for researchers and developers.

\section{Background: Multi-Head Attention and GPU Characteristics}
\label{sec:background}

\subsection{Multi-Head Attention}
\label{subsec:multi_head_attn}

Consider the query, key, and value sequences for a single head, denoted by $\mathbf{Q}, \mathbf{K}, \mathbf{V} \in \mathbb{R}^{N \times d}$, where $N$ represents the length of the sequence and $d$ the dimensionality of each head. The resulting attention output $\mathbf{O}$ can be determined by:

\begin{equation*}
  \mathbf{S} = \alpha \mathbf{Q} \mathbf{K}^\top \in \mathbb{R}^{N \times N}, \quad \mathbf{P} = \mathrm{softmax}(\mathbf{S}) \in \mathbb{R}^{N \times N}, \quad \mathbf{O} = \mathbf{P}\mathbf{V} \in \mathbb{R}^{N \times d},
\end{equation*}

In this context, the $\mathrm{softmax}$ function operates across the rows, and the scaling constant is commonly chosen as $\alpha = 1/\sqrt{d}$. To maintain numerical stability during exponentiation, it is standard to subtract $\mathrm{rowmax}(\mathbf{S})$ from $\mathbf{S}$. For multi-head attention (MHA), every attention head utilizes distinct projection matrices for the queries, keys, and values. These calculations are executed in parallel for all heads and across different batches, ultimately yielding the complete output tensor.

Let $\phi$ denote a scalar loss function, and use $\mathbf{d}(-) = \partial \phi / \partial (-)$ to represent the gradient operator. Assuming we have the output gradient $\mathbf{dO} \in \mathbb{R}^{N \times d}$, we apply the chain rule to determine $\mathbf{dQ}$, $\mathbf{dK}$, and $\mathbf{dV}$ as described below:

\begin{align*}
  \mathbf{dV} &= \mathbf{P}^\top \mathbf{dO} \in \mathbb{R}^{N \times d} \\
  \mathbf{dP} &= \mathbf{dO} \mathbf{V}^\top \in \mathbb{R}^{N \times N} \\
  \mathbf{dS} &= \mathrm{dsoftmax} (\mathbf{dP}) \in \mathbb{R}^{N \times N} \\
  \mathbf{dQ} &= \alpha \mathbf{dS} \mathbf{K} \in \mathbb{R}^{N \times d} \\
  \mathbf{dK} &= \alpha \mathbf{dS}^\top \mathbf{Q} \in \mathbb{R}^{N \times d},
\end{align*}

In this context, we use the relation $\mathbf{d}s = (\mathrm{diag}(p) - p p^\top)\mathbf{d}p$, where $p = \mathrm{softmax}(s)$ and $s$ is a vector. The notation $\mathrm{dsoftmax}(\mathbf{dP})$ denotes the application of this expression to each row independently. Additionally, this operation can be efficiently performed in parallel over the set of heads and batches during the backward propagation of MHA.

\subsection{GPU hardware characteristics and execution model}
\label{subsec:hardware}

We discuss elements of the GPU execution model that are pertinent to \textsc{FlashAttention-3}, emphasizing the NVIDIA Hopper architecture as a specific example of this framework.

\paragraph{Memory hierarchy:} GPU memory is structured into a layered hierarchy, where higher bandwidth is associated with smaller capacities (\cref{tab:gpu-hierarchy})\footnote{\citet{luo2024benchmarking} describes shared memory bandwidth as 128 bytes per clock cycle per SM, which is multiplied by 132 SMs and a boost clock rate of 1830 MHz for calculation.}. The most extensive memory, referred to as global memory (GMEM) or HBM, consists of off-chip DRAM that is shared among all streaming multiprocessors (SMs). When data is fetched from GMEM, it is automatically cached into an on-chip L2 cache. Beyond this, each SM features its own dedicated, highly banked shared memory (SMEM), which is a small, programmer-controlled on-chip cache. At the lowest level of the hierarchy lies the register file within every SM.

\paragraph{Thread hierarchy:} In the GPU programming paradigm, threads are structured into various logical groupings that facilitate execution. Starting from the smallest unit, the hierarchy consists of individual threads, then warps, which contain 32 threads each, followed by warpgroups composed of 4 consecutive warps. These are aggregated into threadblocks (also known as cooperative thread arrays or CTAs), which may be further grouped into threadblock clusters (as seen in Hopper architectures), and ultimately all are encompassed within grids.

The two hierarchies exhibit a strong interconnection. Threads belonging to a given CTA are simultaneously assigned to the same SM, while CTAs grouped within a cluster are concurrently assigned to the same GPC. SMEM is accessible to every thread inside its CTA, in contrast to RMEM, where each thread is allotted no more than 256 private registers.

\paragraph{Asynchrony and warp-specialization:}

GPUs function as throughput-oriented processors that leverage parallelism and asynchronous operations to mask both execution and memory delays. In Hopper, the Tensor Memory Accelerator (TMA) serves as specialized hardware to facilitate asynchronous memory transfers between GMEM and SMEM \cite[\S7.29]{cuda}. Moreover, differing from previous architectures like Ampere, Hopper’s Tensor Core—accessible through the warpgroup-scoped WGMMA instruction \cite[\S9.7.14]{ptx}—also operates asynchronously and is capable of receiving its input data directly from shared memory.

Asynchronous hardware capabilities enable the design of warp-specialized kernels, wherein the warps within a CTA are assigned distinct roles as either producers (responsible exclusively for data transfer) or consumers (focused solely on computation). This specialization generally enhances the compiler’s potential to construct more efficient instruction schedules \citep{warp-specialization-2011}. Furthermore, Hopper architecture introduces support for dynamic register reassignment among warpgroups using \verb|setmaxnreg| \cite[\S9.7.17.1]{ptx}, thereby allowing warps performing MMAs to access a greater portion of RMEM, while warps issuing TMA—tasks that require only a single thread—receive fewer resources.

\paragraph{Low-precision number formats:}
\label{sec:low-precision-gpu}
Current GPU architectures incorporate dedicated hardware components designed to expedite low-precision calculations. For instance, utilizing the WGMMA instruction with Hopper's FP8 Tensor Cores enables double the throughput per SM relative to operations performed with FP16 or BF16 precision.

Nevertheless, effectively utilizing FP8 WGMMA requires familiarity with the specific layout requirements imposed on its input operands. Consider the scenario of performing a GEMM operation, multiplying $A \times B^{\top}$, where $A$ is an $M \times K$ matrix and $B$ is an $N \times K$ matrix. An operand—either $A$ or $B$—is described as \emph{mn-major} if its outer $M$ or $N$ dimension is stored contiguously, whereas it is called \emph{k-major} if the inner $K$-dimension is stored contiguously. For FP16 WGMMA, input operands from SMEM can be provided in either mn-major or k-major layouts. In contrast, FP8 WGMMA restricts input operands to the k-major format exclusively. Additionally, scenarios such as attention mechanisms—where back-to-back GEMM operations are combined within a single kernel—face additional challenges, as conflicting layouts between the FP32 accumulator and the FP8 operand may hinder the proper execution of successive, dependent FP8 WGMMAs.

Regarding attention mechanisms, these layout constraints necessitate specific alterations to the FP8 algorithm’s design; further details are presented in \cref{sec:algofp8}.

\subsection{Standard Attention and Flash Attention} As described by \citet{dao2022flashattention}, \textbf{standard attention} refers to the approach wherein the attention mechanism on the GPU explicitly writes the intermediate matrices $\mathbf{S}$ and $\mathbf{P}$ to high bandwidth memory (HBM). \textsc{FlashAttention}, on the other hand, introduces a strategy that applies a localized version of softmax reduction, thereby eliminating the costly intermediate memory transactions and consolidating the entire attention operation into a single kernel call. This localized softmax mechanism is implemented through lines \ref{code-ws:softmax_start}-\ref{code-ws:softmax_end} within the consumer mainloop presented in \cref{alg:flash3_wgmma_ws_only}, and involves rescaling segments of $\mathbf{O}$. A straightforward justification confirming that this method correctly produces $\mathbf{O}$ is presented in \cite[\S 2.3.1]{dao2023flashattention2}.

\section{FlashAttention-3: Algorithm}
\label{sec:algo}

In this section, we present the \textsc{FlashAttention-3} algorithm. To streamline the exposition, our primary focus will be on the forward pass; the approach for the backward pass is outlined in~\cref{sec:algo_ws_bwd}. We begin by outlining the method for incorporating warp-specialization in conjunction with a circular SMEM buffer into the foundational \textsc{FlashAttention-2} algorithm. Next, we detail the utilization of WGMMA asynchrony to construct a two-stage pipeline in which GEMM and softmax operations are overlapped. Lastly, we discuss the necessary adaptations for supporting FP8, addressing both layout requirements and the preservation of accuracy through the use of block quantization and handling of incoherent processing.

\subsection{Producer-Consumer asynchrony through warp-specialization and pingpong scheduling}
\label{sec:algo_ws}

\paragraph{Warp-specialization}
Similar to \textsc{FlashAttention-2}, the forward computation in \textsc{FlashAttention-3} exhibits straightforward parallelism across the batch dimension, the set of attention heads, and the query sequence positions. Therefore, providing a CTA-level abstraction of the procedure suffices, where the focus is on processing a specific tile $\mathbf{Q}_i$ of the query matrix to obtain the corresponding tile $\mathbf{O}_i$ in the output. For clarity, we first present the warp-specialization strategy assuming a circular SMEM buffer and \emph{without} the complication of overlapping GEMM and softmax operations. Denote $d$ as the dimensionality of each head, $N$ as the total length of the sequence, and choose a fixed query tile size $B_r$ that partitions $\mathbf{Q}$ into $T_r = \lceil \frac{N}{B_r} \rceil$ distinct tiles $\mathbf{Q}_1, ..., \mathbf{Q}_{T_r}$.

\begin{algorithm}[H]
    \caption{\small\label{alg:flash3_wgmma_ws_only}\textsc{FlashAttention-3} forward pass \textbf{excluding} intra-consumer overlap -- CTA perspective}
    \begin{algorithmic}[1]
\REQUIRE Matrices $\mathbf{Q}_i \in \mathbb{R}^{B_r \times d}$ and $\mathbf{K}, \mathbf{V} \in \mathbb{R}^{N \times d}$ stored in HBM, key block size $B_c$ where $T_c = \lceil \frac{N}{B_c} \rceil$.
\STATE Set up a pipeline manager to control barrier synchronization using a circular SMEM buffer with $s$ stages.
\IF {executing within a producer warpgroup}
\STATE Release a specified quantity of registers.
\STATE Trigger transfer of $\mathbf{Q}_i$ from HBM to shared memory.
\STATE After load finishes, signal consumer that $\mathbf{Q}_i$ has been fetched.
\FOR{$0 \le j < T_c$}
    \STATE Await the consumption of the buffer's $(j\,\%\,s)$th stage.
    \STATE Start transferring $\mathbf{K}_j, \mathbf{V}_j$ from HBM to shared memory at buffer stage $(j\,\%\,s)$.
    \STATE Once transfer is done, notify consumers regarding the readiness of $\mathbf{K}_j, \mathbf{V}_j$.
\ENDFOR
\ELSE \STATE Allocate a predetermined set of registers, contingent upon the number of consumer warps.
\STATE Initialize on-chip: $\mathbf{O}_i = (0) \in \mathbb{R}^{B_r \times d}$, $\ell_i = 0$ and $m_i = -\infty$ within $\mathbb{R}^{B_r}$.
\STATE Block until $\mathbf{Q}_i$ is available in shared memory.
\FOR{$0 \le j < T_c$}
\STATE Block until $\mathbf{K}_j$ appears in shared memory.
\STATE Execute $\mathbf{S}_i^{(j)} = \mathbf{Q}_i \mathbf{K}_j^T$ via SS-GEMM. Commit, then wait. \label{alg:ws_only_gemm1}
\STATE Save $m_i^{\mathrm{old}} = m_i$; update $m_i = \mathrm{max}(m_i^{\mathrm{old}}, \mathrm{rowmax}(\mathbf{S}_i^{(j)}))$. \label{code-ws:softmax_start}
\STATE Calculate $\widetilde{\mathbf{P}}_i^{(j)} = \mathrm{exp}(\mathbf{S}_i^{(j)} - m_i)$, and update $\ell_i = \mathrm{exp}(m_i^{\mathrm{old}} - m_i)\ell_i + \mathrm{rowsum}(\widetilde{\mathbf{P}}_i^{(j)})$. \label{code-ws:softmax_end}
\STATE Block until $\mathbf{V}_j$ is present in shared memory.
\STATE Update $\mathbf{O}_i = \mathrm{diag}(\mathrm{exp}(m_i^{\mathrm{old}} - m_i))^{-1} \mathbf{O}_i + \widetilde{\mathbf{P}}_i^{(j)} \mathbf{V}_j$ using RS-GEMM. Commit and wait. \label{alg:ws_only_gemm2}
\STATE Permit the producer to reclaim the $(j\,\%\,s)$th buffer stage.
\ENDFOR
\STATE Finalize with $\mathbf{O}_i = \mathrm{diag}(\ell_i)^{-1} \mathbf{O}_i$ and $L_i = m_i + \log(\ell_i)$.
\STATE Store results $\mathbf{O}_i$ and $L_i$ back to HBM as the $i$th segment of $\mathbf{O}$ and $L$.
\ENDIF
\end{algorithmic}
\end{algorithm}

In our realization of \cref{alg:flash3_wgmma_ws_only} targeting Hopper, we employ \verb|setmaxnreg| to manage register (de)allocation. The TMA functionality is used to load $\mathbf{Q}_i$ along with the collection $\{ \mathbf{K}_j, \mathbf{V}_j \}_{0 \leq j < T_c}$. The consumer mainloop utilizes WGMMA for carrying out the GEMM operations, with the prefixes SS and RS specifying whether the primary input comes from shared memory or the register file, respectively. It is important to recognize, for proper comprehension of the execution pattern in \cref{alg:flash3_wgmma_ws_only}, that TMA load operations do not block subsequent loads, thanks to their asynchronous behavior. Additionally, in the producer mainloop, waits are absent during the initial $s$ iterations as the buffer is in the process of being populated.

\paragraph{Pingpong scheduling} The asynchronous execution capabilities of WGMMA and TMA, in conjunction with warp-specialized processing, make it feasible to interleave the softmax operations for one warpgroup while another warpgroup performs GEMM. To illustrate the benefits of such overlapping, consider that non-matmul operations on current hardware accelerators have substantially lower throughput than matmul operations. For instance, on the H100 SXM5 GPU, the peak FP16 matmul performance is 989 TFLOPS, while the throughput for special functions like exponentiation is just 3.9 TFLOPS\footnote{According to the CUDA programming guide, each streaming multiprocessor (SM) can execute 16 special function operations per cycle. With 132 SMs operating at 1830 MHz, this results in a total throughput of 3.9 TFLOPS for special functions.}. In an FP16 attention forward pass with a head dimension of 128, the number of FLOPS dedicated to matmul is 512 times greater than that for exponentials, yet the hardware processes exponentials at a rate that is 256 times lower, meaning that exponential evaluations can account for up to 50\% of the runtime per cycle relative to matmul. This bottleneck intensifies with FP8 arithmetic, as matmul throughput increases further, but the throughput for exponentials remains unchanged.

Because the exponential operation is handled by a dedicated hardware component (the multi-function unit), optimal utilization would entail scheduling the exponential computation concurrently with matmul operations on the Tensor Cores. To achieve this, we employ synchronization barriers (\texttt{bar.sync} instructions) to enforce an ordering such that the GEMMs (GEMM1 -- $\mathbf{P} \mathbf{V}$ of one iteration and GEMM0 -- $\mathbf{Q} \mathbf{K}^\top$ of the subsequent iteration) executed by warpgroup 1 precede those of warpgroup 2. Consequently, during the execution of the GEMMs in warpgroup 2, the softmax kernel for warpgroup 1 is dispatched, creating an overlap. This pattern is then reversed: warpgroup 2 computes softmax while warpgroup 1 processes GEMMs, following a ``pingpong'' scheduling pattern, as shown in~\cref{fig:pingpong_scheduling}. While the real-world implementation of pingpong scheduling is less orderly than the diagram suggests, we observe notable performance benefits in practice (for example, FP16 forward pass performance rises from 570 TFLOPS to between 620 and 640 TFLOPS with a head dimension of 128 and sequence length of 8192).

\paragraph{Attention variants} In the case of multi-query attention~\citep{shazeer2019fast} and grouped query attention~\citep{ainslie2023gqa}, we employ the strategy described in \textsc{FlashAttention-2}, modifying the tensor indices to prevent redundancy of $\mathbf{K}$ and $\mathbf{V}$ within HBM.

\subsection{Intra-warpgroup overlapping GEMMs and softmax}

Within a single warpgroup, it is possible to interleave certain operations from the softmax computation with specific steps in the GEMMs. We outline a method for achieving this overlap.

Within the attention mechanism, computations inside the primary loop exhibit sequential dependencies that restrict opportunities for intra-iteration parallelism. For instance, the (local) softmax calculation (from lines \ref{code-ws:softmax_start} to \ref{code-ws:softmax_end}) is contingent on the output $\mathbf{S}_i^{(j)}$ generated by the initial GEMM operation. Likewise, the subsequent GEMM operation requires the output $\widetilde{\mathbf{P}}_i^{(j)}$ from the softmax stage as an input. In practice, wait statements at lines \ref{alg:ws_only_gemm1} and \ref{alg:ws_only_gemm2} in \cref{alg:flash3_wgmma_ws_only} enforce serialization between softmax and the GEMM operations. To alleviate these dependencies, one can employ inter-iteration pipelining by utilizing supplementary register-based buffers. Building on this strategy, we introduce a pipelined algorithm for GEMM-softmax with two stages,\footnote{It should be noted that the number of overlapping pipeline stages is limited by, though not necessarily equal to, the number $s$ of stages allocated in the circular SMEM buffer.} described as follows:

\begin{algorithm}[H]
    \caption{\small\label{alg:flash3_wgmma}\textsc{FlashAttention-3} consumer warpgroup forward pass}
    \begin{algorithmic}[1]
      \REQUIRE  Matrices $\mathbf{Q}_i \in \mathbb{R}^{B_r \times d}$ and $\mathbf{K}, \mathbf{V} \in \mathbb{R}^{N \times d}$ stored in HBM, key block size $B_c$, where $T_c = \lceil \frac{N}{B_c} \rceil$.
      \STATE Adjust the allocation of registers according to the count of consumer warps.
      \STATE On the chip, set initial values: $\mathbf{O}_i = (0) \in \mathbb{R}^{B_r \times d}$, with $\ell_i = (0)$ and $m_i = (-\infty)$, both in $\mathbb{R}^{B_r}$.
      \STATE Ensure both $\mathbf{Q}_i$ and the initial key block $\mathbf{K}_0$ have been loaded into shared memory.
      \STATE With WGMMA, calculate $\mathbf{S}_{\mathrm{cur}} = \mathbf{Q}_i \mathbf{K}_0^T$; then commit and await completion.
      \STATE Free the buffer stage assigned to the $0$th segment of $\mathbf{K}$.
      \STATE Update $m_{i}$, compute $\tilde{\mathbf{P}}_{\mathrm{cur}}$, and determine $\ell_{i}$ from $\mathbf{S}_{\mathrm{cur}}$, applying rescaling to $\mathbf{O}_i$ as necessary.
      \FOR{$1 \le j < T_c - 1$}
        \STATE Wait until $\mathbf{K}_j$ resides in shared memory.
        \label{code:mainloop_start}
        \STATE Using WGMMA, evaluate $\mathbf{S}_{\mathrm{next}} = \mathbf{Q}_i \mathbf{K}_{j}^T$; commit the operation but do not synchronize.
        \label{code:first_wgmma}
        \STATE Pause until $\mathbf{V}_{j-1}$ has been transferred to shared memory.
        \STATE With WGMMA, execute $\mathbf{O}_{i} = \mathbf{O}_{i} + \tilde{\mathbf{P}}_{\mathrm{cur}} \mathbf{V}_{j-1}$; commit without waiting.
        \label{code:second_wgmma}
        \STATE Wait for WGMMA's computation of $\mathbf{Q}_i \mathbf{K}_{j}^T$ to finalize.
        \STATE Derive $m_{i}$, $\tilde{\mathbf{P}}_{\mathrm{next}}$, and $\ell_{i}$ based on the results in $\mathbf{S}_{\mathrm{next}}$.
        \label{code:softmax}
        \STATE Wait until $\tilde{\mathbf{P}}_{\mathrm{cur}} \mathbf{V}_{j-1}$ WGMMA has finished, then rescale $\mathbf{O}_i$.
        \STATE Release the $(j\,\%\,s)$th buffer slot for $\mathbf{K}$ and $(j-1\,\%\,s)$th for $\mathbf{V}$.
        \STATE Assign the data from $\mathbf{S}_{\mathrm{next}}$ to $\mathbf{S}_{\mathrm{cur}}$.
        \label{code:mainloop_end}
      \ENDFOR 
      \STATE Wait until $\mathbf{V}_{T_c - 1}$ has been made available in shared memory.
      \STATE Compute $\mathbf{O}_{i} = \mathbf{O}_{i} + \tilde{\mathbf{P}}_{\mathrm{last}} \mathbf{V}_{T_c - 1}$ with WGMMA, and wait for completion.

\STATE Epilogue: Adjust $\mathbf{O}_{i}$ by scaling with $m_{i}$. Determine $L_{i}$ using $m_{i}$ and $\ell_{i}$.  
      Store $\mathbf{O}_{i}$ and $L_{i}$ in HBM as the $i$-th segments of $\mathbf{O}$ and $L$.  
    \end{algorithmic}
  \end{algorithm}

\cref{alg:flash3_wgmma} serves as a substitute for the consumer segment in \cref{alg:flash3_wgmma_ws_only}, thereby providing the full \textsc{FlashAttention-3} implementation for FP16 precision. In this context, we treat WGMMA as synonymous with asynchronous GEMM. During the main loop (spanning lines \ref{code:mainloop_start} to \ref{code:mainloop_end}), the execution of the second WGMMA operation in iteration $j$ (see line \ref{code:second_wgmma}) is scheduled to coincide with the softmax computations from the subsequent iteration $j+1$ (refer to line \ref{code:softmax}).

Although the pipelined architecture discussed previously presents potential for enhanced performance, several real-world factors must be addressed:

\paragraph{Compiler reordering} The pseudocode provides a simplified execution order, but in reality, the compiler (NVCC) may reorder instructions to improve efficiency. Such reordering can interfere with the intended sequence of WGMMA and accompanying operations, potentially leading to deviations from the desired overlap and less-than-expected performance benefits. Examination of the generated SASS code confirms that, in this case, the compiler produces the intended overlapped execution (Section~\ref{sec:2-stage-sass}).

\paragraph{Register pressure} Optimal throughput requires keeping register spilling to a minimum. However, implementing the 2-stage pipeline introduces greater register usage, as additional storage is necessary for intermediate computations and to preserve state across stages. Notably, an extra $\mathbf{S}_{\mathrm{next}}$ variable must reside in registers, resulting in supplemental register consumption of $B_r \times B_c \times \text{sizeof}(\text{float})$ per threadblock. This heightened demand for registers can limit feasible block sizes—another tuning knob that commonly drives up register usage. In practice, careful profiling is essential to strike a balance between these competing demands.

\paragraph{3-stage pipelining} Building on the previously outlined 2-stage pipeline, a 3-stage extension is also proposed to further increase overlap—here, the second WGMMA is pipelined concurrently with the softmax calculation. While this could push Tensor Core utilization higher still, it also calls for an even greater register count due to the added pipeline stage, thus exacerbating the challenge of balancing tile dimension and pipeline depth. An in-depth exposition of the 3-stage methodology and corresponding empirical results is given in ~\cref{sec:3-stage}.

\subsection{Low-precision with FP8}
\label{sec:algofp8}

\textbf{Efficiency: layout transformations.} Executing the forward pass of \textsc{FlashAttention-3} using FP8 precision introduces further complications with respect to layout compatibility that are absent when utilizing FP16.

Initially, it is important to observe that the input tensors $\mathbf{Q}$, $\mathbf{K}$, and $\mathbf{V}$ are usually stored so that the head dimension is contiguous in memory. However, to conform to the k-major layout required by FP8 WGMMA for the second GEMM, the $\mathbf{V}$ tensor—more specifically, the tiles of $\mathbf{V}$ loaded into SMEM—must be contiguous along the sequence length dimension. Since TMA loading operations are unable to alter which dimension is contiguous, we have two primary options: (1) perform a transposition of $\mathbf{V}$ in GMEM as a preprocessing step, or (2) carry out an in-kernel transpose of the tiles after they are loaded into SMEM. To achieve (1), it is possible to either (1a) combine the transposition with the epilogue phase of a previous step like the rotary embedding, or (1b) launch a dedicated preprocessing transpose kernel\footnote{A well-optimized transpose kernel can approach device bandwidth limits~\citep{colfax_cutlass_transpose_2024}.}, which effectively swaps the strides between the sequence length and head dimensions. Nevertheless, integrating (1a) into standard libraries presents significant challenges, while (1b) imposes a prohibitive memory bandwidth overhead, which is especially undesirable during inference when memory resources are at a premium.

For FP8 \textsc{FlashAttention-3}, we select option (2). During the in-kernel transpose step, we utilize the LDSM (\verb|ldmatrix|) and STSM (\verb|stmatrix|) instructions, which enable a group of threads within a warp to efficiently transfer data between SMEM and RMEM in blocks of 128 bytes.\footnote{While the PTX specification characterizes LDSM/STSM as handling $8 \times 8$ matrices composed of 16-bit elements \cite[\S 9.7.13.4.15-16]{ptx}, it is possible to pair 8-bit values together to make use of these instructions in FP8 scenarios. However, since the transpose variants of LDSM/STSM are unable to separate packed 8-bit elements, additional register manipulations are necessary between LDSM and STSM to perform a correct tile-wise transpose; for brevity, further details are omitted.} Employing LDSM and STSM not only conserves registers, which makes it feasible to use them within the producer warpgroup, but also facilitates data layout transpositions as part of the memory copying process. Additionally, after the initial step, we can schedule the transpose of the subsequent $\mathbf{V}$ tile to be carried out concurrently with the execution of the two WGMMAs that process the previous $\mathbf{V}$ tile and the current $\mathbf{K}$ tile.

Second, it is evident that, in contrast to FP16, the FP32 accumulator’s memory arrangement in FP8 WGMMA does not align with the register layout expected for operand A. This distinction is illustrated through examples in \cref{fig:rmem_accum} and \cref{fig:rmem_operand}, where the sequence in which each thread stores register entries is displayed. By utilizing byte permute instructions, it becomes feasible to reformat the accumulator resulting from the first WGMMA into a representation compatible with that needed by the subsequent WGMMA, and consistent with the structure of the $\mathbf{V}$ tile generated during the in-kernel transpose. Concretely, referring to \cref{fig:rmem_accum}, we reorder the data as follows:
$$\{ \verb|d0 d1 d4 d5 d2 d3 d6 d7| \},$$
and this specific register permutation is repeated for every set of 8 bytes. From the perspective of the $\mathbf{P}$ tile’s logical arrangement, this operation effectively permutes the tile’s columns (for example, columns $0189$ are shifted to become the leading four columns). In order for WGMMA to subsequently produce the correct output tile, we can design the in-kernel transpose to emit a corresponding row permutation in the $\mathbf{V}$ tile.\footnote{This flexibility achieved through in-kernel transpose removes the need for shuffle instructions for inter-thread register ownership transfers, as previously discussed in~\citep{colfax_fp8_flashattention_2024}.}

\textbf{Accuracy: block quantization and incoherent processing.} The FP8 (e4m3) format allocates only 3 bits for the mantissa and 4 bits to the exponent, which introduces greater numerical inaccuracies compared to FP16 or BF16. In addition, large-scale models often exhibit the presence of outlier elements~\citep{dettmers2208llm,
  sun2024massive}, whose magnitudes far exceed those of the bulk of the data, thereby posing challenges for effective quantization. The common practice is to apply per-tensor scaling~\citep{micikevicius2022fp8}, whereby a single scaling factor is maintained for each tensor (such as separate scalars for $\mathbf{Q}$, $\mathbf{K}$, and $\mathbf{V}$).
To address the increased numerical error in FP8-based attention computations, we introduce two mitigation strategies:

\begin{enumerate}

\item \textbf{Block quantization}: In this approach, a single scaling factor is maintained per block. For each of $\mathbf{Q}$, $\mathbf{K}$, and $\mathbf{V}$, the tensor is divided into segments of size $B_r \times d$ or $B_c \times d$, with each segment undergoing quantization independently. This quantization procedure can be combined with the operation immediately preceding the attention mechanism (for instance, rotary embedding), without incurring any additional computational overhead, as rotary embedding is primarily limited by memory bandwidth. Given that the \textsc{FlashAttention-3} algorithm is inherently designed for blockwise operations, we are able to rescale each block of $\mathbf{S}$ to reflect the effects of block quantization without adding to the computational cost.

\item \textbf{Incoherent processing}: To mitigate the influence of outliers, we pre-process $\mathbf{Q}$ and $\mathbf{K}$ by multiplying each with a random orthogonal matrix $\mathbf{M}$ prior to FP8 quantization. Owing to the orthogonality of $\mathbf{M}$, with $\mathbf{M} \mathbf{M}^\top = I$, it follows that $(\mathbf{Q} \mathbf{M}) (\mathbf{K} \mathbf{M})^\top = \mathbf{Q} \mathbf{K}^\top$, ensuring that the attention computation remains unchanged by this transformation. This method redistributes outlier values, as each entry in $\mathbf{Q} \mathbf{M}$ or $\mathbf{K} \mathbf{M}$ results from a randomized combination of the entries in $\mathbf{Q}$ or $\mathbf{K}$, thereby lessening quantization errors. Consistent with the methodology described by \citet{chee2024quip} and~\citet{tseng2024quip}, we construct $\mathbf{M}$ as the product of random diagonal $\pm 1$ matrices and a Hadamard matrix, allowing for multiplication in $O(d \log d)$ time, rather than $O(d^2)$, and enabling fusion with the rotary embedding step at no extra computational expense.

\end{enumerate}

We demonstrate in \cref{sec:numerical_error} that these strategies can decrease numerical error by as much as 2.6$\times$.

\section{Empirical Validation}
\label{sec:exp}

Our implementation of \textsc{FlashAttention-3} relies on the WGMMA and TMA abstractions provided by CUTLASS~\citep{Thakkar_CUTLASS_2023}, which we utilize to assess both its performance and precision.

\begin{itemize}

\item \textbf{Attention benchmarking.} We evaluate the execution time of \textsc{FlashAttention-3} over a range of sequence lengths and benchmark its performance against a conventional PyTorch implementation, \textsc{FlashAttention-2}, a Triton-based version of \textsc{FlashAttention-2} leveraging H100-specific instructions, and a vendor-optimized release of \textsc{FlashAttention-2} for H100 GPUs via cuDNN. Our results demonstrate that \textsc{FlashAttention-3} delivers up to 2.0$\times$ speed improvement compared to \textsc{FlashAttention-2}, and is 1.5$\times$ faster relative to the Triton variant. Furthermore, \textsc{FlashAttention-3} attains throughput of up to 740 TFLOPs/s, which corresponds to 75\% of the peak theoretical performance on H100 GPUs.

\item \textbf{Ablation investigation.} We verify that the introduction of warp-specialization alongside the pipelining of GEMM-softmax are the primary algorithmic enhancements that underpin the observed speed gains of \textsc{FlashAttention-3}.

\item \textbf{FP8 attention precision.} Our assessment confirms that incorporating block quantization and incoherent computation substantially mitigates numerical errors in FP8 \textsc{FlashAttention-3} by a factor of 2.6$\times$.

\subsection{Benchmarking Attention}
\label{subsec:benchmark_attn}

We evaluate the execution time of various attention algorithms on an H100 80GB SXM5 GPU across multiple configurations (with or without a causal mask, and using head dimensions of 64 or 128) employing FP16 data types. The findings are presented in~\cref{fig:benchmark_attn_fwd} and~\cref{fig:benchmark_attn_bwd}. These figures reveal that \textsc{FlashAttention-3} achieves approximately 1.5-2.0$\times$ the speed of \textsc{FlashAttention-2} during the forward computation, and 1.5-1.75$\times$ during the backward computation. In relation to conventional attention mechanisms, \textsc{FlashAttention-3} demonstrates acceleration ranging from 3$\times$ up to 16$\times$. Notably, for sequence lengths of 1,000 tokens and greater, \textsc{FlashAttention-3} even outpaces the performance of a vendor-optimized closed-source library (cuDNN) designed specifically for H100 GPUs.

\paragraph{Benchmark settings:} The sequence length is adjusted across 512, 1k, up to 16k, with the batch size selected such that the aggregate token count remains fixed at 16k. The hidden layer is configured to a dimension of 2048, while the head dimension is chosen as 64, 128, or 256, corresponding to 32, 16, or 8 heads, respectively. FLOPs for the forward pass are computed as follows:

\begin{equation*}
  4 \cdot \text{seqlen}^2 \cdot \text{head dimension} \cdot \text{number of heads}.
\end{equation*}

When applying causal masking, this value is halved, reflecting that roughly 50\% of the elements are actually computed. For the backward pass, we estimate the FLOPs by scaling the forward pass FLOPs by a factor of 2.5, given that the forward computation involves 2 matrix multiplications, whereas the backward computation—with recomputation—requires 5 matrix multiplications.

Additionally, we evaluate the forward pass runtime using FP8 under comparable conditions. The outcomes for a head dimension of 256 are presented in ~\cref{fig:benchmark_attn_fp8}, with comprehensive results detailed in \cref{sec:benchmark_attn_fp8_full}.

\subsection{Ablation Study: 2-Stage Pipelining Experiments}
\label{sec:2-stage-pipelining-experiments}

We conduct an ablation study on the two-stage WGMMA-softmax pipelining and warp-specialization mechanisms in the context of non-causal FP16 \textsc{FlashAttention-3}, maintaining fixed values for the parameters $\{\text{batch}, \text{seqlen}, \text{nheads}, \text{hdim}\} = \{ 4, 8448, 16, 128\}$. As shown in~\cref{table:ablation_pipelining}, these algorithmic enhancements—including asynchronous operations utilizing warp-specialization and the overlap of GEMM with softmax—yield a considerable performance gain, increasing throughput from 570 TFLOPs to 661 TFLOPs.

\subsection{Numerical Error Validation}
\label{sec:numerical_error}

Given ongoing concerns regarding the numerical error associated with \textsc{FlashAttention}~\citep{golden2024flash}, we perform a comparative analysis of \textsc{FlashAttention-2}, \textsc{FlashAttention-3}, and a conventional attention implementation, using a FP64 reference as a baseline. To emulate the outlier features and activations commonly encountered in LLMs~\citep{dettmers2208llm, sun2024massive}, we initialize the elements of $\mathbf{Q}, \mathbf{K}, \mathbf{V}$ based on the distribution described below:

\begin{equation*}
  \mathcal{N}(0, 1) + \mathcal{N}(0, 100) \cdot \mathrm{Bernoulli}(0.001).
\end{equation*}

In this setup, all elements follow a normal distribution with mean zero and a standard deviation of 1; however, for 0.1\% of these elements, an additional independent noise component—normally distributed with a standard deviation of 10—is included. The resulting root mean squared error (RMSE) is reported in~\cref{table:numerical_error}. Under the FP16 format, both \textsc{FlashAttention-2} and \textsc{FlashAttention-3} demonstrate RMSE values that are 1.7$\times$ lower than those from the conventional implementation, as the intermediate computations (specifically, the softmax outputs) utilize FP32 precision. The standard attention method using FP8 applies per-tensor scaling, computes matmul accumulations in FP32, and retains intermediate softmax outcomes in FP16. Owing to the benefits of block quantization and noncoherent processing, \textsc{FlashAttention-3} executed in FP8 attains a 2.6$\times$ accuracy improvement over this reference approach.

\section{Dicussion, Limitations, Conclusion}
\label{sec:discussion}

Through \textsc{FlashAttention-3}, we illustrate that leveraging emerging programming paradigms and hardware capabilities—specifically asynchrony and reduced precision—can significantly enhance both the speed and fidelity of attention mechanisms. Our approach yields a 1.5-2.0$\times$ acceleration in attention computations relative to \textsc{FlashAttention-2}, while achieving a 2.6$\times$ reduction in FP8 numerical error when compared with conventional per-tensor quantization. There are, however, certain constraints in our current study that we plan to explore further: improvements in large language model (LLM) inference optimization, the integration of a persistent kernel into the FP8 kernel,\footnote{Our experiments currently employ a persistent kernel and load balancing scheme for FP16 \textsc{FlashAttention-3}, but not for the FP8 implementation. This contributes to the comparatively lower performance of FP8 \textsc{FlashAttention-3} for smaller sequence lengths and when causal masking is present, relative to FP8 cuDNN kernels.} and a more comprehensive investigation into the implications of low-precision attention during extensive training runs. Although the present work is centered on Hopper GPUs, we anticipate these methodologies will generalize to additional hardware accelerators. We envision that improved efficiency and accuracy in attention primitives can pave the way for novel use cases in processing tasks with longer contexts.

\end{document}